\chapter{Reference Implementation: Singleflight Coalescing}
\label{ch:reference_singleflight}

This appendix provides the reference Python implementation for the Singleflight Request Coalescing architecture \cite{colakel2026singleflight} discussed in Chapter 8. This exact code is utilized within the AMPS \texttt{asyncio} event loop to prevent GPU cache stampedes when hundreds of agents simultaneously request inference for identical macroeconomic shocks.

\section{The Singleflight Mutex Map}

The implementation relies on an \texttt{asyncio.Lock} to guarantee thread-safe modifications to a global dictionary of active Promises (\texttt{asyncio.Future}).

\begin{lstlisting}[language=Python]
import asyncio
from typing import Any, Callable, Coroutine, Dict, Tuple

class Singleflight:
    """
    Prevents cache stampedes by coalescing simultaneous identical async calls 
    into a single GPU inference execution.
    """
    def __init__(self):
        # Maps a unique prompt hash to an active Future
        self._active_calls: Dict[str, asyncio.Future] = {}
        # Tracks the number of waiters per key
        self._waiter_counts: Dict[str, int] = {}
        # Mutex to protect the dictionary during await blocks
        self._lock = asyncio.Lock()

    async def execute(self, key: str, coro_fn: Callable[[], Coroutine[Any, Any, Any]]) -> Tuple[Any, bool, int]:
        """
        Execute a coroutine. If multiple calls share the same key, only the first 
        call executes. Subsequent calls await the first call's Future.
        
        Returns:
            Tuple containing:
            - The result of the coroutine
            - Boolean indicating if this specific caller executed the coroutine
            - Integer count of how many callers shared this result
        """
        async with self._lock:
            if key in self._active_calls:
                future = self._active_calls[key]
                self._waiter_counts[key] += 1
                is_first = False
            else:
                future = asyncio.get_event_loop().create_future()
                self._waiter_counts[key] = 1
                self._active_calls[key] = future
                is_first = True

        if is_first:
            try:
                result = await coro_fn()
                future.set_result(result)
                waiters = self._waiter_counts[key]
                return result, True, waiters
            except Exception as e:
                future.set_exception(e)
                raise
            finally:
                async with self._lock:
                    if key in self._active_calls:
                        del self._active_calls[key]
                    if key in self._waiter_counts:
                        del self._waiter_counts[key]
        else:
            result = await future
            # Note: _waiter_counts[key] is deleted by the first caller, so we must rely on a localized snapshot if we want precise numbers for late callers, but for illustration we return a flag.
            return result, False, 0
\end{lstlisting}

\section{Execution Metrics and GPU VRAM Preservation}

By wrapping the \texttt{llama.cpp} inference calls inside this \texttt{Singleflight.execute()} method, AMPS guarantees that when the 671B God Agent broadcasts an 8-K filing, the 400 micro-agents first attempt to execute the \textit{Base Semantic Parsing} prompt. Because this parsing prompt is identical across all agents, the string is hashed to a \texttt{key}. 

Agent 1 acquires the lock, registers the Future, and sends the prompt to the GPU (consuming VRAM for the KV Cache). Agents 2 through 400 instantly hit the `key in self.\_active\_calls` branch and go to sleep. They do not execute their unique, memory-intensive psychological profiles until the base semantic parsing is coalesced and distributed. 

Without this 60-line block of code, Agents 2 through 400 would simultaneously allocate redundant KV caches for the exact same 8-K parsing phase, instantly crashing the simulation.
